\documentclass[12pt]{article}

\usepackage{amsmath,amssymb,amsthm}
\usepackage{tikz-cd}
\usepackage{tikz}
\usepackage{hyperref}
\usepackage{cleveref}
\usepackage{graphicx}
\usepackage[margin=1in]{geometry}
\usepackage{everypage}
\usepackage{xcolor}
\usepackage{booktabs}
\usepackage{enumitem}
\usepackage{mathtools}
\usepackage{microtype}

\hypersetup{
  colorlinks=true,
  linkcolor=blue!55!black,
  citecolor=green!45!black,
  urlcolor=blue!65!black
}

\newtheorem{theorem}{Theorem}[section]
\newtheorem{proposition}[theorem]{Proposition}
\newtheorem{lemma}[theorem]{Lemma}
\newtheorem{corollary}[theorem]{Corollary}
\theoremstyle{definition}
\newtheorem{definition}[theorem]{Definition}
\newtheorem{example}[theorem]{Example}
\newtheorem{assumption}[theorem]{Assumption}
\theoremstyle{remark}
\newtheorem{remark}[theorem]{Remark}
\newtheorem{warning}[theorem]{Warning}

\newcommand{\A}{\mathcal A}
\newcommand{\Ham}{\mathfrak{Ham}}
\newcommand{\Gap}{\mathfrak{Gap}}
\newcommand{\spec}{\operatorname{spec}}
\newcommand{\dist}{\operatorname{dist}}
\newcommand{\supp}{\operatorname{supp}}
\newcommand{\id}{\operatorname{id}}
\newcommand{\status}[1]{\textsc{#1}}

\definecolor{grokgray}{RGB}{110,110,110}
\AddEverypageHook{%
  \ifnum\value{page}=1
    \begin{tikzpicture}[remember picture,overlay]
      \node[
        rotate=90,
        anchor=south,
        font=\Large\sffamily\bfseries\color{grokgray},
        inner sep=0pt
      ] at ([xshift=38pt,yshift=0.52\paperheight]current page.south west)
      {GrokRxiv:2026.07.thermodynamic-gaps\quad[\,math-ph\,]\quad14 Jul 2026};
    \end{tikzpicture}
  \fi
}

\title{Thermodynamic Spectral Gaps in Parameterized Quantum Lattice Systems:\\
Definitions, Stability, and a Condensed Moduli Formulation}

\author{Matthew Long\\
\textit{The YonedaAI Collaboration}\\
\textit{YonedaAI Research Collective}\\
Chicago, IL\\
\texttt{matthew@yonedaai.com} $\cdot$ \url{https://yonedaai.com}}

\date{14 July 2026}

\begin{document}

\maketitle

\begin{abstract}
We isolate the analytic content required to speak about a uniformly gapped
family of quantum lattice systems. Three notions that are often compressed
into one phrase are kept separate: the literal finite-volume ground-state
gap, the finite-volume gap above a selected low-energy band, and the bulk gap
of the positive generator in a ground-state GNS representation. We prove
elementary separation and compactness results, give explicit pointwise-gapped
families whose gap infimum vanishes, and formulate a gap-witness object that
records the state, exhaustion, boundary convention, spectral band, and common
lower bound. We then state the exact role of Lieb-Robinson estimates,
quasi-adiabatic spectral flow, Local Topological Quantum Order, and controlled
perturbation theorems. Condensed and profinite parameter objects organize
families and descent, but do not supply a spectral gap or its stability. Our
main proposal separates a fibration of quantitative gap witnesses from its
qualitative uniformly gapped image subfunctor. The construction is conditional on independent
analytic input and is designed to prevent pointwise gappedness from being
mistaken for a thermodynamic phase. We close with precise theorem interfaces
for formal libraries and with limitations relevant to microscopic and
effective classifications of topological matter.
\end{abstract}

\tableofcontents

\section{Introduction}

A spectral gap is one of the principal analytic inputs in the mathematical
definition of a zero-temperature quantum phase. It controls adiabatic
transport, protects low-energy information, supports clustering statements,
and allows topological invariants to remain locally constant along a family.
It is also a frequent source of hidden assumptions. A finite matrix has a
positive difference between distinct eigenvalues, but that observation says
nothing about a lower bound uniform in system size. A family can have a
positive gap at every parameter while its infimum is zero. A bulk GNS
Hamiltonian can be gapped while open boundaries carry gapless modes.

These distinctions become unavoidable in a condensed-mathematics formulation
of phases. A condensed or profinite test object can encode a family of
interactions. Sheaf descent can glue compatible local descriptions. Neither
operation proves that the thermodynamic generator has an isolated ground
sector. The gap must remain a witnessed analytic property.

This paper addresses the third analytic problem in a five-paper program:

\[
\begin{tikzcd}[row sep=large]
\text{local interactions}
  \arrow[d]
\\
\text{parameterized Hamiltonian object}
  \arrow[d]
\\
\text{uniformly gapped subobject}
  \arrow[d]
\\
\text{stable phase groupoid}
  \arrow[d]
\\
\text{invertible phase spectrum.}
\end{tikzcd}
\]

The paper does not construct the final spectrum. It supplies a defensible
meaning for the third line and explains which established stability results
can be imported.

\subsection{Contributions}

Our contributions are the following.

\begin{enumerate}[label=\arabic*.]
  \item We give separate definitions for literal finite-volume gaps,
  ground-band gaps, uniform finite-volume gaps, GNS bulk gaps, and uniform
  parameterized GNS gaps.

  \item We prove an elementary compactness theorem for a constant-rank
  isolated spectral cluster and show why rank change invalidates the common
  compactness shortcut.

  \item We give three explicit counterexamples: a norm-continuous compact
  matrix family with pointwise but nonuniform literal gaps, a locally
  continuous profinite spin family with uniform locality but collapsing GNS
  gaps, and a sequence of finite ferromagnetic chains whose gaps close in the
  thermodynamic limit.

  \item We formulate a fibration of quantitative gap witnesses and its
  propositional image inside the Hamiltonian functor. The witness records
  enough data to make the predicate meaningful and stable under restriction
  of the parameter object.

  \item We separate two established analytic mechanisms. Spectral flow
  transports a low-energy sector along a path already known to have a uniform
  isolated band. Perturbative gap stability preserves a pre-existing gap only
  under substantive structural hypotheses such as frustration freeness,
  quantitative local gaps, and LTQO.

  \item We state formalization contracts for Lean and executable finite-model
  checks for Haskell. These artifacts represent hypotheses and finite
  calculations. They do not claim proofs of infinite-volume theorems.
\end{enumerate}

\subsection{Status vocabulary}

We use five labels throughout:

\begin{center}
\begin{tabular}{ll}
\toprule
Label & Meaning \\
\midrule
\status{established} & A cited theorem or standard construction under stated assumptions.\\
\status{proposed} & A definition or moduli construction introduced here.\\
\status{conjectural} & A precise assertion not proved here.\\
\status{open} & No proof or counterexample is known in the stated scope.\\
\status{obstructed} & A theorem or explicit counterexample rules out the claim.\\
\bottomrule
\end{tabular}
\end{center}

The distinction is important. Formal vocabulary should expose uncertainty,
not conceal it.

\section{Quantum lattice framework}

\subsection{Local and quasi-local observables}

Let $(\Gamma,d)$ be a countable metric space. We assume uniform local
finiteness when a volume-growth estimate is required. Each site $x\in\Gamma$
carries a finite-dimensional Hilbert space $\mathcal H_x$. For a finite set
$\Lambda\Subset\Gamma$, define

\[
\A_\Lambda
=
\bigotimes_{x\in\Lambda}\mathcal B(\mathcal H_x).
\]

If $\Lambda\subset\Lambda'$, the canonical inclusion is

\[
\iota_{\Lambda,\Lambda'}(A)=A\otimes 1_{\Lambda'\setminus\Lambda}.
\]

These maps are unital and isometric. The local algebra and its C*-completion
are

\[
\A_{\mathrm{loc}}
=
\bigcup_{\Lambda\Subset\Gamma}\A_\Lambda,
\qquad
\A_\Gamma
=
\overline{\A_{\mathrm{loc}}}^{\|\cdot\|}.
\]

\begin{definition}[Interaction]
An interaction is a map
\[
\Phi:\{X\Subset\Gamma\}\longrightarrow\A_{\mathrm{loc}}
\]
such that $\Phi(X)=\Phi(X)^*\in\A_X$. Its finite-volume Hamiltonian on
$\Lambda$ is
\[
H^\Phi_\Lambda=\sum_{X\subseteq\Lambda}\Phi(X).
\]
\end{definition}

The formal infinite sum of all interaction terms is generally not an element
of $\A_\Gamma$. Under decay hypotheses, the thermodynamic object is instead a
strongly continuous automorphism group, often defined first on local
observables and then extended.

\subsection{F-functions and uniform locality}

\begin{definition}[F-function]
A nonincreasing function $F:[0,\infty)\to(0,\infty)$ is an F-function on
$(\Gamma,d)$ if
\[
\|F\|
=
\sup_{x\in\Gamma}\sum_{y\in\Gamma}F(d(x,y))
<\infty
\]
and
\[
C_F
=
\sup_{x,y\in\Gamma}
\frac{1}{F(d(x,y))}
\sum_{z\in\Gamma}F(d(x,z))F(d(z,y))
<\infty.
\]
For thermodynamic limits on a general countable metric space, we also require
the uniform tail condition
\[
T_F(R):=
\sup_{x\in\Gamma}
\sum_{d(x,y)\geq R}F(d(x,y))
,\qquad
\lim_{R\to\infty}T_F(R)=0.
\]
This condition is automatic for the standard translation-invariant lattice
examples used below. On $\mathbb Z^\nu$, translation invariance makes the sum
independent of the base point $x$, so it is the ordinary tail of a convergent
series. We state the condition separately because pointwise summability is not
a substitute for a family-uniform tail on arbitrary geometry.
\end{definition}

For an interaction $\Phi$, set

\[
\|\Phi\|_F
=
\sup_{x,y\in\Gamma}
\frac{1}{F(d(x,y))}
\sum_{X\ni x,y}\|\Phi(X)\|.
\]

Standard choices on $\mathbb Z^\nu$ are a sufficiently fast polynomial tail
and its exponentially weighted version,

\[
F(r)=\frac{1}{(1+r)^{\nu+\epsilon}},
\qquad
F_a(r)=e^{-ar}F(r).
\]

\begin{definition}[Uniformly local parameter family]
Let $S$ be a topological parameter space. A family
$\Phi=(\Phi_s)_{s\in S}$ is uniformly F-local if
\[
M_F(\Phi)=\sup_{s\in S}\|\Phi_s\|_F<\infty.
\]
It is locally F-continuous if its restrictions to every fixed finite region
vary continuously in the corresponding restricted interaction norm.
\end{definition}

Local F-continuity is a finite-region condition. It does not include the
separate geometric hypothesis $T_F(R)\to0$, which will be stated whenever a
family-uniform thermodynamic limit is used.

Uniform locality and parameter continuity are distinct. Global norm
continuity $S\to\mathcal B_F$ on compact $S$ implies uniform boundedness. It
is often too strong for product-topology disorder because a change far from a
fixed observation region can remain large in a supremum over all sites.

\subsection{Uniform Lieb-Robinson control}

The standard Lieb-Robinson argument yields a bound whose constants depend on
the F-function, the geometry, and $\|\Phi\|_F$. A uniform family bound therefore
gives constants independent of $s$.

\begin{theorem}[Uniform family Lieb-Robinson estimate]
\label{thm:uniform-lr}
Assume $M_F(\Phi)<\infty$. Let $A\in\A_X$ and $B\in\A_Y$ have finite disjoint
supports. For finite $\Lambda\supset X\cup Y$,
\[
\|[\tau_t^{\Lambda,s}(A),B]\|
\leq
\frac{2\|A\|\|B\|}{C_F}
\left(e^{2C_FM_F(\Phi)|t|}-1\right)
\sum_{x\in X}\sum_{y\in Y}F(d(x,y)).
\]
The same right side works for every $s\in S$.
\end{theorem}

\begin{proof}[Proof status]
This is an \status{established} uniform-family corollary of the usual
interaction-norm Lieb-Robinson theorem. The proof is the standard iterated
commutator estimate. Every occurrence of the interaction norm is bounded by
the common constant $M_F(\Phi)$. We do not claim a new proof of the sharp
boundary form of the estimate.
\end{proof}

\begin{corollary}[Uniform thermodynamic dynamics]
\label{cor:thermo-dynamics}
Under the assumptions of \cref{thm:uniform-lr}, and with the additional
uniform-tail hypothesis $T_F(R)\to0$, finite-volume dynamics converge on local
observables along any regular exhaustion. The convergence is uniform in $s$
and on compact time intervals. If the family is also locally F-continuous, the
map $s\mapsto\tau_t^s(A)$ is norm-continuous for every local $A$.
\end{corollary}

\begin{proof}[Proof sketch]
Compare dynamics in two nested volumes by a Duhamel expansion. The difference
is bounded by a common F-tail outside the smaller region, multiplied by the
uniform exponential factor from \cref{thm:uniform-lr}. Summability of $F$
and the separately assumed limit $T_F(R)\to0$ make this tail vanish uniformly
in the base point. The same estimate separates a parameter variation
inside a large finite region from the common tail outside it.
\end{proof}

\begin{warning}
Neither \cref{thm:uniform-lr} nor \cref{cor:thermo-dynamics} contains a gap
hypothesis or a gap conclusion. Uniform causal control and a uniform spectral
gap are logically separate properties.
\end{warning}

\section{Five gap predicates}

\subsection{Literal finite-volume ground-state gap}

Let $H_\Lambda$ be self-adjoint on the finite-dimensional Hilbert space for
$\Lambda$. Write

\[
E_0(\Lambda)=\min\spec(H_\Lambda).
\]

\begin{definition}[Literal ground-state gap]
The literal finite-volume ground-state gap is
\[
\gamma_{\mathrm{gs}}(\Lambda)
=
\inf\bigl(\spec(H_\Lambda)\setminus\{E_0(\Lambda)\}\bigr)
-E_0(\Lambda),
\]
with the convention that every eigenvector at exactly $E_0(\Lambda)$ belongs
to the ground space.
\end{definition}

This definition is sensitive to a low-energy level that approaches the ground
energy and becomes exactly degenerate in a limit.

\subsection{Finite-volume ground-band gap}

\begin{definition}[Selected ground-band gap]
Choose a compact interval $I_\Lambda$ whose intersection with the spectrum is
the selected low-energy cluster, and assume that the complementary spectral
set is nonempty. Define
\[
\gamma_{\mathrm{band}}(\Lambda)
=
\dist\left(
\spec(H_\Lambda)\cap I_\Lambda,
\spec(H_\Lambda)\setminus I_\Lambda
\right).
\]
The rank of the selected spectral projection is part of the convention.
\end{definition}

This is the natural finite-volume predicate for spectral flow. An isolated
low-energy band can survive small internal splittings even when the literal
ground-state gap becomes small.

\subsection{Uniform finite-volume gap}

\begin{definition}[Uniform finite-volume band gap]
Fix an exhaustion $(\Lambda_n)$, a boundary convention, and selected bands
$I_{s,n}$. A family has uniform finite-volume gap $\gamma>0$ if
\[
\inf_{s\in S}\inf_n
\gamma_{\mathrm{band}}(s,\Lambda_n)
\geq\gamma.
\]
\end{definition}

The exhaustion and boundary convention cannot be suppressed. Open and
periodic volumes can display different low-energy behavior.

\subsection{GNS bulk gap}

Let $\omega$ be an infinite-volume ground state for the dynamics $\tau$. Let
$(\mathcal H_\omega,\pi_\omega,\Omega_\omega)$ be its GNS triple. The dynamics
is implemented by a positive generator $H_\omega$ with
$H_\omega\Omega_\omega=0$ \cite{BratteliRobinson2}.

\begin{definition}[GNS bulk gap]
The GNS gap of the selected ground state is
\[
\gamma_{\mathrm{GNS}}(\omega)
=
\sup\{\gamma>0:(0,\gamma)\cap\spec(H_\omega)=\varnothing\}.
\]
\end{definition}

If the GNS ground vector is unique, the positive gap can be expressed through
a Poincare-type estimate on the domain of the generator. The representation
and state remain essential data.

\subsection{Uniform parameterized GNS gap}

\begin{definition}[Uniform bulk-gapped family]
A family of interactions with selected ground states $(\omega_s)_{s\in S}$ is
uniformly bulk-gapped if there is $\gamma>0$ such that
\[
\gamma_{\mathrm{GNS}}(\omega_s)\geq\gamma
\qquad
\text{for all }s\in S.
\]
\end{definition}

Pointwise positivity of $\gamma_{\mathrm{GNS}}(\omega_s)$ is not enough.
Uniformity is a separate quantifier and is the physically relevant one for a
single parameterized phase object.

\section{Compactness: valid theorem and invalid shortcut}

Compactness does produce a common separation for a finite-dimensional spectral
cluster when its rank is fixed. The rank condition is the missing hypothesis
in a common but invalid argument.

\begin{proposition}[Compact constant-rank spectral separation]
\label{prop:constant-rank}
Let $S$ be compact and let $H:S\to M_N(\mathbb C)_{\mathrm{sa}}$ be
norm-continuous. Suppose there is a continuous family of rank-$r$ spectral
projections $P_s$ and that
\[
\spec(H_s|_{P_s\mathbb C^N})
\cap
\spec(H_s|_{(1-P_s)\mathbb C^N})
=\varnothing
\]
for every $s$. Then the distance between the two spectral clusters has a
positive minimum on $S$.
\end{proposition}

\begin{proof}
The ordered eigenvalues of a self-adjoint matrix are Lipschitz continuous in
the operator norm. After choosing which $r$ eigenvalues form the cluster, the
distance between the two finite sets of continuous eigenvalue functions is a
continuous positive function on $S$. A positive continuous function on a
compact space has a positive minimum.
\end{proof}

\begin{example}[Rank change destroys the conclusion]
\label{ex:matrix-collapse}
Let
\[
S=\{0\}\cup\{1/n:n\geq2\}
\]
and
\[
H_s=\operatorname{diag}(0,s,1).
\]
This is a norm-continuous family. At $s=1/n$, the literal ground-state gap is
$1/n$. At $s=0$, the ground space has rank two and the literal gap is $1$.
Every fiber is literally gapped, but
\[
\inf_{s\in S}\gamma_{\mathrm{gs}}(H_s)=0.
\]
If the two lowest levels are selected as one band, the band gap is uniformly
positive. The disagreement is not a paradox. The two predicates encode
different low-energy sectors.
\end{example}

\begin{corollary}
Compactness plus pointwise literal gappedness does not imply a uniform literal
gap when the ground-state rank changes.
\end{corollary}

\section{Thermodynamic counterexamples}

\subsection{A moving weak defect}

Let $\Gamma=\mathbb N$ with one qubit at each site. Let
$p_x=|1\rangle\langle1|_x$. Consider the product interactions whose formal
Hamiltonians are

\[
H_\infty=\sum_{x\geq1}p_x,
\qquad
H_n=\frac{1}{n}p_n+\sum_{x\neq n}p_x.
\]

The parameter space $S=\mathbb N\cup\{\infty\}$ is profinite. On every fixed
finite region, $H_n$ agrees with $H_\infty$ for all sufficiently large $n$.
The family is locally continuous and uniformly finite range. Its
Lieb-Robinson constants can therefore be chosen uniformly.

\begin{proposition}[Uniform locality without a uniform gap]
Each system above has the same unique product ground state. Its GNS excitation
gap is
\[
\gamma_{\mathrm{GNS}}(H_n)=1/n,
\qquad
\gamma_{\mathrm{GNS}}(H_\infty)=1.
\]
Hence
\[
\inf_{s\in S}\gamma_{\mathrm{GNS}}(H_s)=0.
\]
\end{proposition}

\begin{proof}
A single spin flip at site $n$ has energy $1/n$ for $H_n$, and every nonzero
excitation has energy at least $1/n$. For $H_\infty$, a single spin flip has
energy $1$ and all other excitations have integer energy at least $1$.
\end{proof}

The family is not continuous in the global interaction norm that takes a
supremum over all sites. This example explains why local parameter continuity
and uniform locality should be stated separately.

\subsection{Finite chains with a closing gap}

For a length-$L$ open spin-$1/2$ ferromagnetic chain, let

\[
H_L
=
\sum_{x=1}^{L-1}\frac{1-P_{x,x+1}}{2},
\]

where $P_{x,x+1}$ swaps neighboring spins. Each summand is positive and the
model is frustration free.

\begin{proposition}[Finite positivity does not imply a thermodynamic gap]
Every finite $H_L$ has a positive gap above its symmetric ground space, but
\[
0<\gamma_L
\leq
1-\cos(\pi/L)
\longrightarrow0.
\]
\end{proposition}

\begin{proof}[Proof sketch]
Restrict to the one-magnon sector. There the Hamiltonian is proportional to
the graph Laplacian of the length-$L$ path. Its first nonzero eigenvalue is
$1-\cos(\pi/L)$ in the stated normalization. The variational principle gives
the upper bound on the full many-body gap.
\end{proof}

The infinite isotropic ferromagnetic chain is gapless. A theorem about all
finite volumes must contain a lower bound independent of $L$ if it is to imply
a thermodynamic gap.

\section{Spectral flow transports an assumed gap}

Quasi-adiabatic continuation relates uniformly gapped paths to quasi-local
automorphisms. It is a transport theorem, not a gap-existence theorem.

\subsection{Hypotheses}

Let $s\mapsto\Phi(s)$ be differentiable for $s\in[0,1]$. The standard
finite-volume-to-thermodynamic spectral-flow theorem assumes, in a typical
form:

\begin{enumerate}[label=(\alph*)]
  \item finite-dimensional local spin spaces;
  \item a regular exhaustion $(\Lambda_n)$ with controlled geometry;
  \item a decay norm for $\Phi(s)$ and $\partial_s\Phi(s)$ bounded uniformly
  in $s$;
  \item a selected finite-volume spectral interval $I_s$;
  \item a separation $\gamma>0$ between that interval and the rest of the
  spectrum, uniform in $s$ and $n$.
\end{enumerate}

Choose a filter $W_\gamma$ adapted to the gap. The finite-volume generator is

\[
D_{\Lambda_n}(s)
=
\int_{\mathbb R}W_\gamma(t)
\tau_t^{H_{\Lambda_n}(s)}
\bigl(\partial_sH_{\Lambda_n}(s)\bigr)
\,dt.
\]

The decay of the filter and the Lieb-Robinson estimate make this generator
quasi-local.

\begin{theorem}[Thermodynamic spectral flow]
\label{thm:spectral-flow}
Under the stated locality, differentiability, geometry, and uniform
ground-band gap assumptions, the finite-volume flows converge to a strongly
continuous cocycle $\alpha_s$ of quasi-local automorphisms of $\A_\Gamma$.
The thermodynamic ground-state sets satisfy
\[
\mathcal S(s)=\mathcal S(0)\circ\alpha_s.
\]
\end{theorem}

\begin{proof}[Proof status]
This is the \status{established} automorphic-equivalence theorem of Bachmann,
Michalakis, Nachtergaele, and Sims, stated here in the notation of this paper.
The proof combines a filtered quasi-adiabatic generator, quasi-locality bounds,
uniform convergence, and transport of the isolated spectral projection.
\end{proof}

\begin{warning}[No circular use]
The uniform gap is an input to the filter and to the spectral projection
transport. \Cref{thm:spectral-flow} cannot be cited to prove that the same path
has a gap.
\end{warning}

\subsection{Bulk spectral-flow variants}

There are bulk formulations based on a unique selected ground state with a
uniform GNS gap and suitable differentiability estimates. These avoid treating
open-boundary low-energy modes as bulk gap closure. They still assume the
bulk gap. They do not establish general perturbative stability.

\section{Controlled perturbative gap stability}

General stability of the spectral gap for arbitrary interacting Hamiltonians
is not available. Strong theorems exist for structured classes. We describe a
representative frustration-free lane.

\subsection{Frustration-free reference system}

Assume an interaction indexed by sites,

\[
H_0=\sum_{x\in\Gamma}h_x,
\]

where $h_x\geq0$ is supported in a ball $b_x(R)$ of common finite radius,
$\sup_x\|h_x\|<\infty$, and the selected infinite-volume state has zero
energy on every local term.

Let $H_0^{\mathrm{GNS}}$ be the positive GNS generator. A stability theorem
starts with

\[
\gamma_0=\operatorname{gap}(H_0^{\mathrm{GNS}})>0.
\]

This is a pre-existing bulk gap, not a consequence of frustration freeness.

\subsection{Local gap control}

For controlled local regions $\Lambda(x,n)$, assume a polynomial lower bound
on the local gap,

\[
\spec(H_{\Lambda(x,n)})
\subset
\{0\}\cup[\gamma_1n^{-\alpha},\infty)
\]

for constants $\gamma_1>0$ and $\alpha\geq0$. The local gap may decay, but its
decay is quantified.

\subsection{Local Topological Quantum Order}

Let $P_{b_x(m)}$ project onto the local ground space in a ball of radius $m$.
A representative LTQO estimate is

\[
\left\|
P_{b_x(m)}AP_{b_x(m)}
-\omega_0(A)P_{b_x(m)}
\right\|
\leq
\|A\|(1+k)^\nu G_0(m-k)
\]

for $A$ supported in $b_x(k)$ and $m\geq k$. The decay function $G_0$ must
have enough finite moments relative to the spatial growth, partition growth,
and local-gap exponent.

This estimate says that local ground states become indistinguishable by an
observable sufficiently far from the boundary. Its quantitative form is
essential to the perturbative argument.

\subsection{Anchored perturbations}

Write the perturbation as an anchored interaction

\[
V=\sum_{x\in\Gamma}\sum_{n\geq0}\Phi(x,n),
\qquad
\Phi(x,n)\in\A_{b_x(n)}^{\mathrm{sa}},
\]

with stretched-exponential decay

\[
\|\Phi(x,n)\|
\leq
\|\Phi\|e^{-an^\theta},
\qquad
a>0,
\quad
0<\theta\leq1.
\]

\begin{theorem}[Controlled bulk gap stability]
\label{thm:controlled-stability}
Assume regular polynomial geometry, a finite-range positive
frustration-free reference interaction, a selected GNS sector whose vacuum
vector is nondegenerate for the implementing bulk generator,
a positive bulk GNS gap $\gamma_0$, quantitative local gap control, LTQO with
the required decay moments, and a stretched-exponentially decaying anchored
perturbation. For every $0<\gamma<\gamma_0$, there is a perturbation threshold
$s_0(\gamma)>0$ such that $|s|<s_0(\gamma)$ preserves a unique transported
ground state and a GNS gap at least $\gamma$.
\end{theorem}

\begin{proof}[Proof status]
This is an \status{established} theorem in the Nachtergaele-Sims-Young
stability framework. The proof uses spectral flow to conjugate the perturbed
Hamiltonian into a form whose ground-state contribution is controlled by
LTQO, while the local gap bounds dominate the remaining relatively bounded
terms. We state the hypothesis package because omitting any part would suggest
a false theorem for arbitrary gapped systems.
\end{proof}

\begin{remark}[Finite-volume topological degeneracy]
Nondegeneracy in \cref{thm:controlled-stability} refers to the vacuum vector in
the selected infinite-volume GNS sector. It does not assert that the local
ground projections $P_{b_x(m)}$ have rank one. LTQO is designed to control
locally indistinguishable finite-volume ground spaces, including topologically
degenerate ones.
\end{remark}

\begin{remark}
Earlier finite-volume stability results of Michalakis and collaborators
provide a closely related LTQO lane. Precise assumptions differ between
formulations. A paper must cite the version it actually uses.
\end{remark}

\subsection{What the theorem does not prove}

\Cref{thm:controlled-stability} does not apply automatically to:

\begin{itemize}
  \item arbitrary frustrated Hamiltonians;
  \item power-law interactions outside the allowed decay class;
  \item reference systems without a known bulk gap;
  \item models lacking local gap control or LTQO;
  \item families in which the symmetry or ground-state convention changes;
  \item claims about all open-boundary spectra.
\end{itemize}

\section{Witness fibration and qualitative gapped subfunctor}

We now state the moduli proposal. It is designed so that gap data cannot be
lost by notation.

\subsection{Hamiltonian family object}

Let $S$ be a compact Hausdorff or profinite test object. Define schematically

\[
\Ham_F(S)
=
\left\{
(\Phi_s)_{s\in S}:
\sup_s\|\Phi_s\|_F<\infty,
\text{ with local F-continuity}
\right\}.
\]

Additional data records symmetry, on-site Hilbert spaces, and the allowed
coarse geometry. The assignment should be treated as a groupoid or higher
stack once unitary equivalences and automorphisms are included.

\subsection{Gap witness}

\begin{definition}[Gap witness]
A gap witness for $\Phi\in\Ham_F(S)$ consists of:

\begin{enumerate}[label=(\roman*)]
  \item a declared gap predicate, either finite-volume band or GNS bulk;
  \item an exhaustion and boundary convention in the finite-volume case;
  \item a selected spectral band of locally constant rank, or selected ground
  states $\omega_s$ in the GNS case;
  \item a common number $\gamma>0$;
  \item verification that the selected predicate is at least $\gamma$ for all
  $s$ and all required volumes;
  \item continuity data for the selected projections or states sufficient for
  the later phase equivalence.
\end{enumerate}
\end{definition}

Define the quantitative witness object

\[
\Gap_F^{\mathrm{wit}}(S)
=
\left\{
(\Phi,\mathsf w):
\Phi\in\Ham_F(S),
\ \mathsf w\text{ is a gap witness}
\right\}.
\]

The projection
\[
p_S:\Gap_F^{\mathrm{wit}}(S)\longrightarrow\Ham_F(S)
\]
forgets the witness. It is not generally injective because one family can have
many lower bounds, states, exhaustions, or theorem providers. Thus the
witnessed object is a fibration of data over the Hamiltonian functor, not a
subfunctor in the monomorphism sense.

Define the qualitative gapped image by propositional existence:
\[
\Gap_F^{\mathrm{prop}}(S)
=
\{\Phi\in\Ham_F(S):\|\exists\mathsf w,\
(\Phi,\mathsf w)\in\Gap_F^{\mathrm{wit}}(S)\|\}.
\]
Here $\|{-}\|$ means that only existence is retained. Equivalently, take the
image of $p_S$. The qualitative object is a subfunctor of $\Ham_F$ when image
formation and pullback are implemented in the chosen sheaf or stack category.
Both constructions are \status{proposed}.

\begin{proposition}[Restriction stability]
Let $u:T\to S$ be continuous. Pullback of a witnessed family along $u$
preserves the common lower bound and defines a witnessed family over $T$.
\end{proposition}

\begin{proof}
The pulled-back family is indexed by $t\mapsto u(t)$. The uniform locality
bound over $T$ is no larger than the bound over $S$. Every selected gap is one
of the gaps already bounded below on $S$, so the same $\gamma$ works. Pull back
the state, band, and continuity data componentwise.
\end{proof}

\subsection{Descent is not gap creation}

Suppose a finite jointly surjective family $S_i\to S$ covers the parameter
object. Compatible interaction families can descend. If every local witness
comes with the same explicit lower bound and compatible state or band data,
the witness can also descend in a controlled formulation.

It is invalid to assume only that every restriction has some unspecified
positive lower bound. The minimum is available for a finite cover, but the
selected bands and states must agree on overlaps. For varying or infinite
covers, further uniformity is required. Sheaf descent organizes compatible
evidence. It does not produce evidence that was not supplied.

If gapped neighborhoods form an open cover of a compact parameter space, then
compactness provides a finite subcover. Compatible quantitative witnesses on
that finite subcover yield a common positive minimum. This argument still
requires one fixed predicate and compatible state or band conventions. A bare
pointwise statement need not provide such neighborhoods or compatibility.

\subsection{Openness and discriminant locus}

For a norm-continuous finite matrix family with a fixed-rank isolated cluster,
gappedness is open. For interacting thermodynamic families, openness follows
only inside a class covered by a stability theorem.

Given a parameter map $f:B\to\Ham_F$, define the witnessed gapped locus

\[
U_f=B\times_{\Ham_F}\Gap_F^{\mathrm{prop}}
\]

and the gapless discriminant

\[
\Sigma_f=B\setminus U_f.
\]

On underlying parameter sets or spaces, this pullback is a genuine subobject,
so its complement is meaningful. The larger pullback with
$\Gap_F^{\mathrm{wit}}$ is instead a bundle of witnesses over $U_f$ and must
not be subtracted from $B$. The notation is meaningful only after the gap
predicate, image construction, and witness category are fixed. A topological
phase label is expected to be locally constant on $U_f$ after quotienting by
the chosen stable gapped equivalences.

\section{Relations among the predicates}

The following table summarizes safe implications.

\begin{center}
\begin{tabular}{p{0.32\textwidth}p{0.56\textwidth}}
\toprule
Statement & Status and qualification \\
\midrule
Uniform finite-volume gap implies a bulk GNS gap
& \status{established} in controlled limits with stated uniqueness,
regularity, and boundary assumptions. Not a definition-level implication.\\[3pt]
Bulk GNS gap implies open-boundary finite-volume gap
& \status{obstructed} in general by topological boundary modes.\\[3pt]
Every finite volume is gapped implies a thermodynamic gap
& \status{obstructed} by the ferromagnetic-chain example.\\[3pt]
Pointwise parameter gaps imply a common gap
& \status{obstructed} without constant-rank or stability hypotheses.\\[3pt]
Uniform locality implies a uniform gap
& \status{obstructed} by the moving weak-defect example.\\[3pt]
Uniform gapped path implies quasi-local automorphic transport
& \status{established} under spectral-flow locality and regularity assumptions.\\[3pt]
Small perturbations preserve an arbitrary interacting gap
& \status{open} as a universal theorem; \status{established} only for
controlled structural classes.\\
\bottomrule
\end{tabular}
\end{center}

\section{Worked finite examples}

\subsection{Two-level constant-rank family}

Let

\[
H_s=
\begin{pmatrix}
-1 & s\\
s & 1
\end{pmatrix},
\qquad
s\in[-1,1].
\]

The eigenvalues are $\pm\sqrt{1+s^2}$, so the separation is
$2\sqrt{1+s^2}\geq2$. The negative spectral projection has constant rank one.
This is the elementary setting where compactness and continuity correctly
produce a common gap.

\subsection{A stable product-spin perturbation}

For $N$ qubits, consider

\[
H_0=\sum_{x=1}^N p_x,
\qquad
V=\sum_{x=1}^N v_x,
\]

with $\|v_x\|\leq\epsilon$ and commuting diagonal $v_x$. If every excited
on-site level remains at least $1-2\epsilon$ above the local ground level,
then the many-body product gap is at least $1-2\epsilon$, independent of $N$.
This direct estimate uses the product and commuting structure. It is not a
general theorem for noncommuting perturbations.

\subsection{Bulk and boundary conventions}

An open topological chain can have boundary modes with a splitting that
vanishes exponentially in length, while the periodic bulk has a stable gap.
The literal open-chain ground-state gap, the open-chain ground-band gap, and
the bulk GNS gap answer different questions. A phase definition intended to
ignore protected edge degeneracy should use a band or bulk predicate and say
so explicitly.

\section{Formal and executable representations}

\subsection{Lean theorem interfaces}

A Lean library can honestly represent the logical dependency graph without
formalizing the full analytic proofs. It should contain:

\begin{itemize}
  \item finite sites, finite supports, and finite-volume Hamiltonian data;
  \item a parameter family and a predicate expressing a common norm bound;
  \item literal finite-spectrum gaps and selected-band gaps;
  \item a structure containing a positive common gap witness;
  \item an abstract GNS-gap predicate whose provider is an explicit interface;
  \item an abstract spectral-flow provider that requires a uniform gap witness;
  \item an abstract stability provider that lists LTQO and local-gap inputs;
  \item status and provenance fields for each imported analytic result.
\end{itemize}

The preferred design uses structures and theorem-provider typeclasses rather
than global axioms. A concrete development can instantiate the interface when
the relevant theorem has been formalized. No declaration should suggest that
Lean proved an infinite-volume spectral gap merely because a finite list of
eigenvalues was checked.

\subsection{Haskell demonstrations}

Executable code accompanying this paper checks finite calculations:

\begin{enumerate}[label=\arabic*.]
  \item ordered-spectrum gap extraction;
  \item selected-band separation;
  \item the matrix family in \cref{ex:matrix-collapse};
  \item the $L^{-2}$ closing upper bound for the ferromagnetic chain;
  \item finite prefixes of the moving weak-defect gap sequence;
  \item predicates distinguishing pointwise positivity from a declared common
  lower bound.
\end{enumerate}

These checks are \status{computational}. They catch convention and
implementation errors. They are not thermodynamic proofs.

\section{Main results in concise form}

We collect the paper's conclusions.

\begin{theorem}[Separation theorem]
Uniform F-locality, local parameter continuity, and pointwise positive GNS
gaps do not imply a parameter-uniform GNS gap, even for a profinite family of
finite-range commuting product interactions with a common unique ground
state.
\end{theorem}

\begin{proof}
The moving weak-defect family supplies all stated properties and has gaps
$1/n$ along a sequence converging to the parameter $\infty$.
\end{proof}

\begin{theorem}[Witness necessity]
Any moduli assignment that records only an interaction family and the statement
that each fiber is gapped cannot distinguish the moving weak-defect family from
a uniformly gapped family. A uniform phase subobject must therefore record or
require a common quantitative lower bound, together with the gap convention.
\end{theorem}

\begin{proof}
The fiberwise predicate has the same truth value in both cases. The desired
distinction depends on the infimum across parameters and, in finite volume,
across the exhaustion. This information is not recoverable from a bare list of
positive truth values.
\end{proof}

\begin{theorem}[Conditional phase transport]
For a differentiable uniformly local path equipped with a uniform
finite-volume ground-band witness satisfying the hypotheses of
\cref{thm:spectral-flow}, the endpoints have quasi-locally automorphically
equivalent thermodynamic ground-state sets.
\end{theorem}

\begin{proof}
Apply \cref{thm:spectral-flow}. The word ``conditional'' records that the gap
witness is an input rather than a conclusion.
\end{proof}

\begin{theorem}[Controlled local stability]
Within the frustration-free LTQO class described in
\cref{thm:controlled-stability}, the witnessed bulk-gapped predicate is open
under sufficiently small perturbations in the stated anchored interaction
norm.
\end{theorem}

\begin{proof}
The stability theorem provides a positive perturbation radius and a common
lower bound smaller than the original bulk gap. Restrict the proposed gap
image subfunctor to the theorem's structural class.
\end{proof}

\section{Discussion}

\subsection{Why a numerical gap function is insufficient}

Writing $\operatorname{gap}(H)$ suggests a canonical real number. In the
thermodynamic setting, the notation hides choices. Which finite volumes and
boundaries are used? Is a near-degenerate ground multiplet one band or several
levels? Which infinite-volume state and representation are selected? Is the
claim about a bulk generator or an edge spectrum? A gap witness makes these
choices data rather than prose.

\subsection{Why condensed mathematics still helps}

The negative conclusion would be too strong if it said condensed mathematics
adds nothing. It adds a coherent category for continuous and profinite
families, finite-resolution disorder, symmetry objects, and descent. It can
organize restrictions and gluing of an already established uniform gap
witness. It can also place the gapless locus inside a parameterized moduli
problem. The point is narrower: categorical organization does not replace the
many-body estimate.

\subsection{Relation to topological phase transitions}

On a witnessed gapped region $U$, a stable phase invariant should be locally
constant. A path connecting different labels must leave $U$ and meet the
discriminant $\Sigma$. The local topology of $\Sigma$ may carry a relative
charge when a generalized cohomology invariant is defined. This reasoning
depends on the gap predicate. If pointwise but nonuniform families are admitted
as single gapped objects, the phase label can fail to have the stability needed
for this interpretation.

\subsection{Relation to effective field theory}

A low-energy field theory presupposes a controlled separation of scales. A
uniform thermodynamic gap is one possible input for a topological effective
description, but it does not alone construct the continuum approximation.
Conversely, a field-theoretic bordism class does not prove that a microscopic
Hamiltonian realizes a uniform gap. The microscopic-to-effective comparison
must preserve the gap claim as a separate verified property.

\subsection{Noninvertible phases}

The gap predicates developed here apply to invertible and noninvertible
systems. A later invertible phase spectrum captures only the sector with a
stacking inverse. General topological order requires excitation and defect
categories in addition to a gap witness. The analytic foundation is shared,
but the classification object is larger.

\section{Limitations and open problems}

The present work has deliberate limits.

\begin{enumerate}[label=\arabic*.]
  \item We do not prove a new universal gap-stability theorem. We expose the
  assumptions of established controlled theorems.

  \item We do not prove descent for a fully constructed higher stack of GNS
  representations. The witness fibration and its qualitative image are
  proposals whose precise categorical targets remain to be built.

  \item We use finite-dimensional on-site spaces in the principal framework.
  Unbounded on-site terms and bosonic systems require domain and topology
  refinements.

  \item We do not treat mobility gaps as spectral gaps. Mobility-gap invariants
  often require Sobolev or localization algebras and different stability
  inputs.

  \item We do not identify a bulk GNS gap with a boundary gap. A framework for
  bulk-boundary pairs should record both predicates and their comparison map.

  \item We do not assert that the set of uniformly gapped interactions is open
  in every natural interaction topology. Openness is established only inside
  the scope of an applicable stability theorem.

  \item We do not infer a thermodynamic theorem from Haskell tests or a Lean
  interface. Formal provenance remains visible.
\end{enumerate}

The most immediate open problem is to construct a category of parameterized
ground-state representations in which uniform gap witnesses satisfy descent
and in which spectral flow is a natural transformation. A second problem is
to compare finite-volume band witnesses with bulk GNS witnesses under minimal
boundary assumptions. A third is to extend the controlled stability lane to
larger frustrated and long-range classes without hiding model-specific input.

\section{Conclusion}

A uniformly gapped family is more than a family whose fibers are each gapped.
It consists of a locality-controlled interaction family, a declared spectral
predicate, a ground-band or state convention, and one lower bound that survives
the thermodynamic and parameter limits. This paper has given that statement a
precise moduli form.

The key logical pattern is:

\[
\boxed{
\begin{array}{c}
\text{uniform locality}
\quad+\quad
\text{explicit gap witness}
\\[4pt]
\Downarrow\ \text{under stated hypotheses}
\\[4pt]
\text{spectral flow or controlled perturbative stability}
\\[4pt]
\Downarrow
\\[4pt]
\text{stable phase data on the witnessed gapped locus.}
\end{array}}
\]

Condensed mathematics can organize this locus across continuous, profinite,
and inverse-limit parameters. The spectral gap itself remains an analytic
theorem. Preserving that division of labor is necessary for a credible
condensed-mathematical theory of topological phases and transitions.

\appendix

\section{Gap-claim audit protocol}

This appendix gives a compact protocol for checking a thermodynamic gap claim
before it enters a phase-classification argument. It is intended for research
papers, formal interfaces, and code-backed examples.

\subsection{Object declaration}

The claim should first identify the microscopic object:

\begin{enumerate}[label=\textbf{O\arabic*.}]
  \item State the metric lattice or coarse space and its volume-growth or tail
  properties.
  \item State the local Hilbert spaces or local observable algebras.
  \item Give the interaction and its decay norm.
  \item State the symmetry action and whether it is fixed along the family.
  \item Distinguish a bounded observable from a formal interaction generator.
\end{enumerate}

Failure at this stage means that no thermodynamic spectral predicate has yet
been defined.

\subsection{Spectral predicate declaration}

The claim should then choose exactly one primary predicate:

\begin{enumerate}[label=\textbf{G\arabic*.}]
  \item Literal finite-volume ground-state gap.
  \item Finite-volume separation above a selected low-energy band.
  \item Uniform finite-volume band gap along a declared exhaustion and boundary
  convention.
  \item Bulk GNS gap for a selected infinite-volume ground state.
  \item Uniform bulk GNS gap for a parameterized family of selected states.
\end{enumerate}

If more than one predicate is used, the comparison theorem must be cited with
its uniqueness, regularity, exhaustion, and boundary assumptions. Similar
notation is not a proof that the predicates coincide.

\subsection{Quantifier audit}

For a family indexed by $S$ and volumes indexed by $n$, write the quantifiers
in full. The desired statement is typically

\[
\exists\gamma>0\ \forall s\in S\ \forall n\geq n_0:
\gamma_{s,n}\geq\gamma.
\]

The weaker statement

\[
\forall s\in S\ \forall n\ \exists\gamma_{s,n}>0:
\gamma_{s,n}\geq\gamma_{s,n}
\]

is nearly automatic for a finite matrix after exact degeneracies are grouped.
It contains no common thermodynamic scale. Changing the order of these
quantifiers changes the physical assertion.

\subsection{Evidence ladder}

Classify the available evidence at one of four levels:

\begin{center}
\begin{tabular}{p{0.18\textwidth}p{0.70\textwidth}}
\toprule
Level & Required evidence \\
\midrule
Finite computation
& Exact diagonalization or certified finite arithmetic for declared sizes. It
does not establish a thermodynamic lower bound.\\[4pt]
Scaling evidence
& A reproducible finite-size study with error bars and an explicit ansatz. It
can support a conjecture but is not a gap theorem.\\[4pt]
Model theorem
& A proof specialized to the interaction, such as a martingale, transfer
operator, integrability, or frustration-free estimate.\\[4pt]
Stability
& Verification that a reference system and perturbation satisfy every
hypothesis of a cited uniform stability result.\\
\bottomrule
\end{tabular}
\end{center}

The result status should match the strongest completed level.

\subsection{Boundary audit}

For each finite-volume statement, record:

\begin{itemize}
  \item open, periodic, twisted, or symmetry-breaking boundary conditions;
  \item whether boundary terms are uniformly bounded and local;
  \item whether edge modes belong to the selected ground band;
  \item whether the claimed phase invariant is bulk, boundary, or relative;
  \item which comparison, if any, relates the finite system to a GNS bulk
  representation.
\end{itemize}

A bulk-boundary correspondence can predict protected edge behavior. It does
not justify calling the open-boundary literal gap a bulk gap.

\subsection{Perturbation audit}

Before citing a stability theorem, verify:

\begin{enumerate}[label=\textbf{P\arabic*.}]
  \item The reference gap is already proved in the theorem's required sense.
  \item The reference interaction belongs to the theorem's structural class.
  \item Local gap or LTQO estimates have the required quantitative decay.
  \item The perturbation belongs to the stated normed interaction class.
  \item The perturbation strength is below a volume-independent threshold.
  \item The selected symmetry and ground-state conventions are preserved.
\end{enumerate}

The phrase ``small local perturbation'' is insufficient without a norm,
threshold, and hypothesis match.

\section{Machine-readable witness schema}

The mathematical witness can be mirrored by a data record. The following is a
language-independent schema rather than an implementation mandate:

\begin{align*}
\mathsf{GapWitness}=\{\,&
\mathsf{predicate},
\mathsf{parameterSpace},
\mathsf{interactionClass},
\mathsf{stateOrBand},\\
&\mathsf{exhaustion},
\mathsf{boundary},
\mathsf{lowerBound},
\mathsf{provider},
\mathsf{status}\,\}.
\end{align*}

Here $\mathsf{provider}$ points to a proof, a cited theorem with discharged
hypotheses, or a computation. The status field prevents a finite numerical
check from being silently promoted to a thermodynamic theorem. Pullback along
a parameter map retains the same witness provider only when every provider
hypothesis is stable under that pullback.

\begin{proposition}[Witness monotonicity]
If a witness proves a lower bound $\gamma>0$, then it also proves every lower
bound $\gamma'$ with $0<\gamma'\leq\gamma$.
\end{proposition}

\begin{proof}
For every indexed spectrum, the inequality
$\operatorname{gap}\geq\gamma$ implies
$\operatorname{gap}\geq\gamma'$ by transitivity of the order on real numbers.
All other witness fields are unchanged.
\end{proof}

\begin{proposition}[Finite-cover minimum]
Suppose a finite cover $S_1,\ldots,S_m$ of $S$ carries compatible witnesses
for the same predicate, state or band convention, exhaustion, and boundary
data, with lower bounds $\gamma_i>0$. Then
\[
\gamma=\min_{1\leq i\leq m}\gamma_i
\]
is a positive global lower bound.
\end{proposition}

\begin{proof}
Every parameter belongs to at least one $S_i$, where its gap is at least
$\gamma_i\geq\gamma$. The minimum of finitely many positive real numbers is
positive. Compatibility is needed to ensure that all local inequalities refer
to the same spectral predicate.
\end{proof}

This elementary proposition explains the useful but limited role of finite
descent. It combines supplied quantitative witnesses. It does not infer them
from qualitative local gappedness.

\begin{thebibliography}{99}

\bibitem{BMNS2012}
S. Bachmann, S. Michalakis, B. Nachtergaele, and R. Sims,
``Automorphic equivalence within gapped phases of quantum lattice systems,''
\emph{Communications in Mathematical Physics} 309 (2012), 835--871.
\url{https://arxiv.org/abs/1102.0842}.

\bibitem{BachmannNachtergaele2014}
S. Bachmann and B. Nachtergaele,
``On gapped phases with a continuous symmetry and boundary operators,''
\emph{Journal of Statistical Physics} 154 (2014), 91--112.
\url{https://arxiv.org/abs/1307.0716}.

\bibitem{BravyiHastingsMichalakis2010}
S. Bravyi, M. Hastings, and S. Michalakis,
``Topological quantum order: stability under local perturbations,''
\emph{Journal of Mathematical Physics} 51 (2010), 093512.
\url{https://arxiv.org/abs/1001.0344}.

\bibitem{BratteliRobinson2}
O. Bratteli and D. W. Robinson,
\emph{Operator Algebras and Quantum Statistical Mechanics 2: Equilibrium
States, Models in Quantum Statistical Mechanics},
second edition, Springer, 1997.

\bibitem{ClausenScholze}
D. Clausen and P. Scholze,
``Lectures on condensed mathematics,'' lecture notes.
\url{https://www.math.uni-bonn.de/people/scholze/Condensed.pdf}.

\bibitem{Hastings2004}
M. Hastings,
``Lieb-Schultz-Mattis in higher dimensions,''
\emph{Physical Review B} 69 (2004), 104431.
\url{https://arxiv.org/abs/cond-mat/0305505}.

\bibitem{KomaNachtergaele1997}
T. Koma and B. Nachtergaele,
``The spectral gap of the ferromagnetic XXZ-chain,''
\emph{Letters in Mathematical Physics} 40 (1997), 1--16.
\url{https://arxiv.org/abs/cond-mat/9512120}.

\bibitem{MichalakisZwolak2013}
S. Michalakis and J. P. Zwolak,
``Stability of frustration-free Hamiltonians,''
\emph{Communications in Mathematical Physics} 322 (2013), 277--302.
\url{https://arxiv.org/abs/1109.1588}.

\bibitem{MoonOgata2020}
A. Moon and Y. Ogata,
``Automorphic equivalence within gapped phases in the bulk,''
\emph{Journal of Functional Analysis} 278 (2020), 108422.
\url{https://arxiv.org/abs/1906.05479}.

\bibitem{NachtergaeleSims2010}
B. Nachtergaele and R. Sims,
``Lieb-Robinson bounds in quantum many-body physics,''
in \emph{Entropy and the Quantum}, Contemporary Mathematics 529 (2010).
\url{https://arxiv.org/abs/1004.2086}.

\bibitem{NSY2019}
B. Nachtergaele, R. Sims, and A. Young,
``Quasi-locality bounds for quantum lattice systems, Part I,''
\emph{Journal of Mathematical Physics} 60 (2019), 061101.
\url{https://arxiv.org/abs/1810.02428}.

\bibitem{NSY2023}
B. Nachtergaele, R. Sims, and A. Young,
``Stability of the bulk gap for frustration-free topologically ordered quantum
lattice systems,''
\emph{Letters in Mathematical Physics} 113 (2023).
\url{https://arxiv.org/abs/2102.07209}.

\bibitem{SimsWarzel}
R. Sims and S. Warzel,
``Decay of determinantal and pfaffian correlation functionals in one-dimensional
lattices,''
\emph{Communications in Mathematical Physics} 347 (2016), 903--931.
\url{https://arxiv.org/abs/1509.00450}.

\end{thebibliography}

\end{document}
